\section{Introduction}
\label{sec:intro}

L-band interferograms are the ones that stay coherent across large earthquakes in vegetated terrain. They are also the
ones the ionosphere distorts most: the dispersive phase scales as $1/f$, and at 1.2\,GHz a modest change in total
electron content makes metre-scale phase ramps and tens of centimetres of curvature across a scene
\citep{meyer2006potential,meyer2011performance}. The split-spectrum method separates the two by their frequency
dependence. The phase of two sub-bands, or of two separate bands, is combined so that the non-dispersive part
(deformation, topography, troposphere) cancels and the dispersive part remains
\citep{rosen2010measurement,brcic2010estimation,gomba2016toward}. It has become the operational correction for
ALOS-2 \citep{liang2018insar,furuya2017midlatitude,nagaoka2025improved}, Sentinel-1 \citep{gomba2017tops,
liang2019ionospheric}, time series \citep{fattahi2017insar,mao2023reformulating,piromthong2026ionospheric} and ice
\citep{liao2018ionospheric}; \citet{mao2025review} review it. NISAR, launched in July 2025, was designed for it.
Its L-band radar transmits a 20, 40 or 77\,MHz main band and a separate 5\,MHz side band 50--70\,MHz above it
\citep{nisar2018handbook}, and the Geocoded Unwrapped Interferogram (GUNW) product carries a split-spectrum
ionospheric phase screen as a layer, to be subtracted by the user \citep{isce3}.

Near a surface rupture the method meets conditions it was not built for. The deformation is steep, discontinuous at
the trace, and often decorrelated over kilometres beside it; large ruptures run to a coast or to the edge of a swath,
cutting the scene into blocks that the unwrapper cannot connect. In principle none of this matters, because the
dispersive combination contains no deformation. In practice, every implementation unwraps something that does:
the two bands separately, or their difference, and then corrects ``unwrapping errors'' against smoothed estimates
\citep{gomba2016toward,fattahi2017insar,isce3}. \citet[Sec.~III-C]{gomba2016toward} already note that their
common-error correction is not robust, because inaccurate estimates are misread as unwrapping errors. At a rupture,
the deformation is exactly such an inaccuracy.

We found the consequence in the NISAR PR GUNWs of two 2026 strike-slip earthquakes: the 24 June events on the San
Sebastián fault in north-central Venezuela, and the 9 October $M_w$\,7.7 earthquake on the Azuero--Son\'a fault zone
in western Panama. Within 10\,km of the ruptures, the GUNW ionospheric screens contain a scaled, sign-flipped copy
of their own track's coseismic deformation. On the four Venezuela tracks the high-passed screen regresses on the high-passed data at $-0.02$ to $-0.12$
($r$ from $-0.19$ to $-0.82$). On the Panama pair it does so at $-0.10$ ($r=-0.81$), and the screen steps by 0.43\,m
of LOS across the trace. Subtracting such a screen adds a few percent of the
deformation back with the opposite sign, which is centimetres of LOS within kilometres of the trace, where slip
inversions read the shallow slip.

This paper does five things.
\begin{enumerate}
  \item It locates where the deformation enters. An ablation runs the two-band, reformulated and reordered
        estimators, with and without ISCE3's correction passes, on identical cells, with an identical filter and
        an identical main-band unwrapping. It uses synthetic ruptures with known truth
        (Section~\ref{sec:synthetic}) and the five NISAR tracks (Section~\ref{sec:nisar}). Two further tests on the
        NISAR tracks feed in the products' own main band and their side-band footprint.
  \item It describes a reordering of the reformulated split-spectrum estimate of \citet{wegmuller2018reformulating}.
        The scaled, unwrapped main band is removed from the wrapped side band \emph{before} anything is unwrapped,
        so the only field that is unwrapped holds no deformation, and no cycle corrections are needed
        (Section~\ref{sec:theory}). The estimate is algebraically the same as the reformulated one, and as the one
        ISCE3 computes; only the order of operations is new.
  \item It shows that the side band ties a block that a cut along the rupture isolates. A wrong main-band integer
        on the block moves the reordered estimate by exactly 3.21\,rad per cycle, and the ionosphere is continuous
        across the gap, so the estimate's step there reads the integer (Section~\ref{sec:sbtie}). On three of the four
        Panama pairs it finds the pipelines' range-offset ties a cycle short. Read the same way, the offsets agree with it.
  \item It states the reordered estimator's remaining error budget near the trace (Section~\ref{sec:nearfield}). On
        synthetics, the cells' handling of steep fringes, the fill's width and the lag of the weighted mean can each
        be removed: by taking the main band out of every pixel, by a fill matched to the gap, and by a weighted local
        plane as the filter's final step. On real data these are millimetre corrections, and next to the D141 rupture
        the main band's local integers set the screen.
  \item It tests the diagnosis on ice (Section~\ref{sec:antarctica}). \citet{joughin2026nisarerrors} found the PR screen
        of two Antarctic pairs offset across a glacier's shear margin. That screen is ours plus whole cycles of the kinds
        above, and where the two differ, the range offsets side with ours. The same boundary is the strong control the
        side-band tie lacked: across a 160\,rad step it ties the cut-off component as the range offsets do.
\end{enumerate}
